\documentclass{article}
\usepackage[margin=1.15in]{geometry}
\usepackage{xcolor}
\usepackage{parskip}
\usepackage{fancyhdr}
\usepackage{array}
\usepackage{booktabs}
\usepackage{enumitem}
\usepackage{amsmath}
\usepackage{tabularx}

\pagestyle{fancy}
\fancyhf{}
\cfoot{\thepage}
\rhead{Lecture 5}
\lhead{MA 214: Applied Statistics (Spring 2026)}
\renewcommand{\headrulewidth}{.4mm}

\newcommand{\blankline}[1]{\underline{\hspace{#1}}}

\begin{document}

\noindent\textbf{Name:}\blankline{2.25in}\hfill\textbf{BUID:}\blankline{1.55in}\newline

\begin{center}
 \Large In-Class Worksheet: Multiple Regression
\end{center}

\subsection*{Scenario}
Multiple regression models a numerical response using more than one predictor. The key new phrase is \textbf{holding the other predictors fixed}. Today we practice coefficient interpretation, indicator variables, and comparing models using adjusted $R^2$.

\medskip
\noindent\textbf{Goals:} \textbf{(1)} read software output and identify reference levels, \textbf{(2)} interpret multiple regression coefficients conditionally, \textbf{(3)} compute predictions from a model with an indicator variable, \textbf{(4)} compare models using $R^2$ and adjusted $R^2$.

\subsection*{Part 1: Reading Software Output}
Interest rates on loans are modeled using \texttt{income\_ver}, a categorical variable with three levels: \texttt{not}, \texttt{source\_only}, and \texttt{verified}. Software reports:

\begin{center}
\renewcommand{\arraystretch}{1.3}
\begin{tabular}{lrrrr}
\toprule
 & \textbf{Estimate} & \textbf{Std. Error} & \textbf{t value} & \textbf{Pr($>|t|$)} \\
\midrule
(Intercept) & 11.0995 & 0.0809 & 137.18 & $<$0.0001 \\
\texttt{income\_ver : source\_only} & 1.4160 & 0.1107 & 12.79 & $<$0.0001 \\
\texttt{income\_ver : verified} & 3.2543 & 0.1297 & 25.09 & $<$0.0001 \\
\bottomrule
\end{tabular}
\end{center}

\begin{enumerate}[label={\bfseries (\Alph*)},itemsep=.8em]
\item The variable has three levels, but only two appear in the table. Why? \\[3em]
\item Which level is the reference level? How can you tell? \\[3em]
\item What is the predicted interest rate for a borrower whose income is \texttt{verified}?
\[
\blankline{3.2in}
\]
\item Which level has the \emph{lowest} predicted interest rate? Explain. \\[3em]
\end{enumerate}

\subsection*{Part 2: Coefficients While Holding Other Variables Fixed}
A data set contains book weights, volumes, and cover types. Define
\[
\text{paperback}=\begin{cases}
1, & \text{paperback book},\\
0, & \text{hardcover book}.
\end{cases}
\]
A fitted model is
\[
\widehat{\text{weight}} = 197.96 + 0.72(\text{volume}) - 184.05(\text{paperback}),
\]
where weight is measured in grams and volume is measured in cm$^3$.

\begin{enumerate}[label={\bfseries (\Alph*)},itemsep=.8em]
\item What is the response variable? What are the predictors? \\[3em]
\item What is the reference level for cover type? \\[2.5em]
\item Interpret the coefficient 0.72 in context. Make sure to include ``holding fixed'' language. \\[4em]
\item Interpret the coefficient $-184.05$ in context. \\[4em]
\item Is the intercept practically meaningful here? Explain briefly. \\[3em]
\end{enumerate}

\subsection*{Part 3: Same Model, Two Lines}
Using the same fitted model,
\[
\widehat{\text{weight}} = 197.96 + 0.72(\text{volume}) - 184.05(\text{paperback}),
\]
complete the two equations below.

\begin{enumerate}[label={\bfseries (\Alph*)},itemsep=.8em]
\item For hardcover books, $\text{paperback}=0$:
\[
\widehat{\text{weight}} = \blankline{4in}
\]
\item For paperback books, $\text{paperback}=1$:
\[
\widehat{\text{weight}} = \blankline{4in}
\]
\item Predict the weight of a hardcover book with volume 900 cm$^3$.
\[
\blankline{3.2in}
\]
\item Predict the weight of a paperback book with volume 900 cm$^3$.
\[
\blankline{3.2in}
\]
\end{enumerate}

\subsection*{Part 4: Comparing Models}
Three models are fit to predict apartment rent.

\begin{center}
\renewcommand{\arraystretch}{1.35}
\begin{tabular}{clcc}
\toprule
\textbf{Model} & \textbf{Predictors} & \boldmath$R^2$ & \textbf{Adjusted }\boldmath$R^2$ \\
\midrule
A & size & 0.61 & 0.60 \\
B & size + bedrooms & 0.64 & 0.62 \\
C & size + bedrooms + random ID code & 0.65 & 0.60 \\
\bottomrule
\end{tabular}
\end{center}

\begin{enumerate}[label={\bfseries (\Alph*)},itemsep=.8em]
\item Which model has the largest $R^2$? \\[2em]
\item Which model has the largest adjusted $R^2$? \\[2em]
\item Which model would you prefer for interpretation? Why? \\[4em]
\item Why is adjusted $R^2$ useful when comparing models with different numbers of predictors? \\[4em]
\end{enumerate}

\subsection*{Wrap-Up}
Complete the sentence: In multiple regression, each coefficient is interpreted as a predicted change in the response \blankline{3.3in}. \\[3em]

\end{document}
